\documentclass[10pt,a4paper]{amsart}
\usepackage[margin=19mm]{geometry}
\usepackage{amsmath,amssymb,mathtools}
\usepackage[scr=rsfs,cal=euler]{mathalfa}
\usepackage{xfrac}
\usepackage[unicode,hidelinks,bookmarks=false]{hyperref}
\newcommand{\ie}{i.e.\ }
\newcommand{\cf}{cf.\ }
\newcommand{\resp}{respectively\ }
\newcommand{\Subsection}{\subsection}
\providecommand{\nobreakdash}{\nobreak\hbox{-}}
\setlength{\emergencystretch}{2em}
\multlinegap0pt
\usepackage{bm}
\usepackage{bbm}
\usepackage{dsfont}
\usepackage{xspace}
\usepackage{cancel}
\usepackage{mathtools}
\usepackage{amsthm}
\newtheorem{theorem}{Theorem}
\newtheorem{lemma}[theorem]{Lemma}
\newtheorem{proposition}[theorem]{Proposition}
\newcommand*{\B}{\mathcal B}% C*-bundles
\newcommand*{\Cst}{\textup C^*}% C*-algebra
\newcommand*{\Id}{\textup{Id}}%identity
\newcommand*{\Locmult}{\mathcal{M}_{\mathrm{loc}}}% local multiplier algebra
\newcommand*{\Null}{\mathcal N}% null space for positive map
\newcommand*{\Slice}{\mathcal S}% inverse semigroup of slices
\newcommand*{\Star}{\(^*\)\nobreakdash-}
\newcommand*{\ess}{\mathrm{e}}% visible instead of essential
\newcommand*{\nb}{\nobreakdash}
\newcommand*{\onto}{\twoheadrightarrow}
\title[C*-irreducible regular inclusions, Galois correspondence and]{C*-irreducible regular inclusions, Galois correspondence and aperiodicity (Proposition 3.4)}
\author{B. K. Kwa\'sniewski, R. Meyer}
\date{Source: arXiv:2609.27518v1 (CC BY 4.0)}
\begin{document}
\maketitle

\noindent\emph{Source and licence.} C*-irreducible regular inclusions, Galois correspondence and aperiodicity. Version arXiv:2609.27518v1 (CC BY 4.0), distributed under the Creative Commons Attribution 4.0 International licence, as recorded in the arXiv licence metadata for this exact version and shown on the abstract page https://arxiv.org/abs/2609.27518v1. Full rights notice: licenses/arxiv-2609.27518-CC-BY-4.0.txt.\par\smallskip

\noindent\textit{Editorial scope.} Counted unit: the Proposition printed as Proposition 3.4 of the article (source0.tex lines 1451--1459, source label \texttt{prop:topological\_grading\_characterisation}), delivered together with the article's own complete original proof of it (source0.tex lines 1461--1506, one \texttt{proof} environment of 46 lines). No mathematics is written, repaired or re-derived: statement and proof are the article's own text line for line. Required context retained uncounted: lem:genuine\_expectation (lines 1367--1377). Every definition, display and label of those lines is retained unchanged. Omitted from this package and not redistributed: 1--1366, 1378--1450, 1460--1460, 1507--4844. Nothing inside a retained excerpt is omitted or altered: every retained body is delivered exactly as the article's lines are written, so no omission receipt of any kind accompanies this package. Citations: the retained excerpts carry no citation command at all, so no bibliography entry of the article is transcribed and no reference is redistributed. Reference adaptation: none was needed and none was made; every reference inside the delivered unit resolves inside it, so no unresolved reference or citation is produced. Printed slips kept literal: no printed word, symbol, hypothesis or number of the counted statement or of its proof was altered. Layout and class: the article's own class is \texttt{amsart} with packages including inputenc, fontenc, lmodern, amssymb, amsxtra, graphicx, mathabx, a4wide, nicefrac, mathtools, amsrefs, enumitem, xy, tikz; the wrapper is the lane's pinned \texttt{amsart} editorial wrapper at 10pt on A4, which restates the theorem-like environments the retained text needs and adds the article's own macro definitions that the retained text needs (\texttt{\textbackslash B}, \texttt{\textbackslash Cst}, \texttt{\textbackslash Id}, \texttt{\textbackslash Locmult}, \texttt{\textbackslash Null}, \texttt{\textbackslash Slice}, \texttt{\textbackslash Star}, \texttt{\textbackslash ess}, \texttt{\textbackslash nb}, \texttt{\textbackslash onto}), with extra wrapper packages (bm, bbm, dsfont, xspace, cancel, mathtools). The delivered numbering is the unit's own. Assets: no retained span contains a figure environment, a graphics-inclusion command, a PDF-page inclusion, a table or a TikZ picture; no image file is copied into this package, the package declares no asset dependency and no image file is redistributed with it. Licence: version arXiv:2609.27518v1 of this article is distributed under the Creative Commons Attribution 4.0 International licence, as recorded in the arXiv licence metadata for this exact version and shown on the abstract page https://arxiv.org/abs/2609.27518v1. A full rights notice is delivered at licenses/arxiv-2609.27518-CC-BY-4.0.txt. Characters: every retained excerpt is pure ASCII, so the delivered engine, XeLaTeX with its default Latin Modern text font, reports no missing character.
\begin{lemma}
  \label{lem:genuine_expectation}
  A \(\Cst\)\nb-inclusion \(A\subseteq B\) with an inverse-semigroup
  grading \(S\subseteq \Slice_A(B)\) admits at most one
  \(S\)\nb-canonical pseudo-expectation~\(E\).
  If it exists, then \(E\colon B\to \Locmult(A)\subseteq I(A)\) takes
  values in the local multiplier algebra.
  Moreover, \(E\) is a genuine expectation that preserves the grading
  in the sense that \(E(M)\subseteq M\) for all \(M\in S\) if and only
  if \((M\cap A)\oplus M^\bot=M\) for all \(M\in S\).
\end{lemma}

\begin{proposition}
  \label{prop:topological_grading_characterisation}
  A \(\Cst\)\nb-inclusion \(A\subseteq B\) is topologically graded if
  and only if~\(B\) is an exotic section \(\Cst\)\nb-algebra for some
  inverse-semigroup action \(\B=(B_{s})_{s\in S}\)  by Hilbert
  bimodules with unit fibre \(B_1=A\).
  If this holds, then \(B\cong\Cst_\ess(\B)\) if and only if  the
  canonical pseudo-expectation for \(A\subseteq B\) is faithful.
\end{proposition}

\begin{proof}
  Let~\(B\) be an exotic section \(\Cst\)\nb-algebra for a saturated
  Fell bundle over a unital inverse semigroup \(\B=(B_{s})_{s\in S}\).
  Let~\(\Phi\) be the quotient map \(\Cst(\B) \onto B\) and
  let~\(E_B\) be the composite of the quotient map \(B\onto
  \Cst_\ess(\B)\) and the faithful pseudo-expectation \(E_\ess\colon
  \Cst_\ess(\B)\to I(A)\).
  We claim that \(E_B\circ \Phi\) is the canonical expectation
  \(E\colon \Cst(\B)\onto A\).
  If \(t\in S\), then \(B_t\cap A\subseteq \Phi(B_t)\cap A\) because
  \(\Phi|_A=\Id_A\).
  Thus \((\Phi(B_t)\cap A)^\bot \subseteq (B_t\cap A)^\bot\) and so
  \[
    E_B(\Phi(B_t)(\Phi(B_t)\cap A)^\bot)
    \subseteq E_B\bigl(\Phi(B_t(B_t\cap A)^\bot)\bigr)
    = E(B_t(B_t\cap A)^\bot)=0.
  \]
  Accordingly, \(E_B\circ\Phi\) is the canonical expectation on
  \(\Cst(\B)\).
  Since~\(E_B\) descends to a faithful map, it is symmetric.
  Hence \(A\subseteq B\) is topologically graded.

  Now let \(A\subseteq B\) be any topologically graded regular
  inclusion.
  In particular, this defines a surjective \Star{}homomorphism
  \(\Phi\colon\Cst(\B) \onto B\).
  Being topologically graded means that the canonical expectation
  \(\Cst(\B)\to I(A)\) descends to a pseudo-expectation on~\(B\).
  The latter must be the unique canonical pseudo-expectation by
  Lemma~\ref{lem:genuine_expectation}.
  Then \(\Phi(\Null_E)\subseteq \Null_{E_B}\) and there is a
  well-defined surjective \Star{}homomorphism
  \[
    \Cst_\ess(\B)= \Cst(\B)/\Null_E \to B/\Null_{E_B},
    \qquad
    b+\Null_E\mapsto \Phi(b)+\Null_{E_B}.
  \]
  This map is injective as well as it intertwines the
  pseudo-expectations \(\Cst(\B)/\Null_E\to I(A)\) and
  \(B/\Null_{E_B}\to I(A)\), which are faithful because both  \(E\)
  and~\(E_B\) are symmetric.
  Composing the quotient map \(B\onto B/\Null_{E_B}\) with
  \(B/\Null_{E_B}\cong \Cst_\ess(\B)\) gives the surjective
  \Star{}homomorphism \(B\onto \Cst_\ess(\B)\) that witnesses
  that~\(B\) is an exotic section \(\Cst\)\nb-algebra of~\(\B\).
\end{proof}

\end{document}
